\begin{figure}[t]
\centering
\resizebox{\linewidth}{!}{%
\begin{tikzpicture}[
  cell/.style={draw=slate!45, dashed, rounded corners=1pt, minimum width=32mm, minimum height=17mm, align=center, font=\scriptsize, fill=mist!35},
  rowlab/.style={anchor=east, font=\scriptsize, align=right},
  colhead/.style={font=\scriptsize\bfseries, align=center}
]
\node[colhead] at (0,1.1) {原图};
\node[colhead] at (3.6,1.1) {金标准};
\node[colhead] at (7.2,1.1) {预测叠加};
\node[colhead] at (10.8,1.1) {误差图};

\foreach \r/\name/\yy in {0/MGD-Former/0, 1/UNet/-2.0, 2/Attention U-Net/-4.0, 3/WRCANet/-6.0, 4/MGD-PriorNet/-8.0, 5/失败案例/-10.0} {
  \node[rowlab] at (-1.4,\yy+0.35) {\name};
  \foreach \c/\x in {0/0,1/3.6,2/7.2,3/10.8} {
    \node[cell] at (\x,\yy) {};
  }
}
\node[font=\scriptsize, color=slate, align=center] at (5.4,-11.2) {占位：真实预测到位后替换为影像；误差图标注假阳性（红）与假阴性（蓝）。};
\end{tikzpicture}}
\caption{定性结果布局（占位）。行为 MGD-Former、四条基线与失败案例；列为原图、金标准、预测叠加与误差图。失败案例需覆盖断裂、粘连、睑缘泄漏与低对比度段漏检。}
\label{fig:qualitative}
\end{figure}
